\documentclass[a4paper,11pt]{article}

% Fonts and symbols
\usepackage{lmodern}
\usepackage{amsmath,amssymb}
\usepackage[utf8]{inputenc}
\usepackage[T1]{fontenc}
\usepackage{inconsolata}
\usepackage{wasysym}
\usepackage{microtype}

% Layout, page setup
\usepackage[margin=2.5cm,a4paper]{geometry}
\usepackage{titlesec}
\usepackage{appendix}
\usepackage{multicol}
\usepackage{multirow}
\usepackage{parskip}
\usepackage{tabto}
\usepackage{pdflscape}
\usepackage{enumitem}
\usepackage{adjustbox}
\usepackage[super]{nth}

% Science and language typesetting
\usepackage{bm}
\usepackage{physics}
\usepackage[
  arc-separator = \,,
  retain-explicit-plus,
  inter-unit-product = \cdot,
  detect-weight=true,
  detect-family=true,
  range-phrase=--,
  range-units=single
]{siunitx}
\usepackage[version=4]{mhchem}
\usepackage{nicematrix}
\usepackage[makeroom]{cancel}
\usepackage{circledsteps}
\usepackage[italian]{babel}
\usepackage{csquotes}

% Graphics, captions, tables
\usepackage[table,dvipsnames]{xcolor}
\usepackage{graphicx}
\usepackage{float}
\usepackage{tikz,tikz-3dplot}
\usepackage{pgfplots}
\pgfplotsset{compat=1.18}
\usepackage{pdfpages}
\usepackage[justification=centering, labelfont={small,bf}, font={small}]{caption}
\usepackage{subcaption}
\usepackage{array}
\usepackage{tabularx}
\usepackage{booktabs}
\usepackage[outline]{contour}

% TikZ options setup
\usetikzlibrary{
  calc,
  arrows,
  arrows.meta,
  positioning,
  decorations.pathreplacing,
  decorations.markings,
  decorations.text,
  calligraphy
}
\usepgfplotslibrary{dateplot}

% References and links
\usepackage{hyperref}
\usepackage[noabbrev]{cleveref}
\usepackage[nottoc,numbib]{tocbibind}
\usepackage{biblatex}

% Miscellaneous
\usepackage{datetime2}
\usepackage{iftex,etoolbox,expl3,xparse}


\title{Architettura BLE lato Client e Server}
\author{Andrea Zorzi -- BillyApp}
\date{Versione 1.1.1 \\ \today}

\begin{document}

\maketitle

\tableofcontents

\newpage

\section{Scopo e perimetro}

Questo documento descrive l’architettura logica e tecnica del sistema BLE lato client e lato server

L'idea è avere un'unica base di riferimento che permetta di:
\begin{itemize}
  \item capire come il client genera, ruota e trasmette l'identificativo BLE \(B_{\text{id}}\);
  \item capire come il backend risolve \(B_{\text{id}}\) in un utente reale;
  \item avere una vista completa sulle scelte di privacy e sicurezza;
  \item alternative valutate e perché sono state scartate o mitigate.
\end{itemize}

Il perimetro è architettuale, non di implementazione. Non sono quindi presenti dettagli di codice, ma solo di logica e flussi.
L'implementazione sarà trattata in un documento a parte.

\section{Panoramica generale}
\subsection{Attori Principali}
Il sistema coinvolge i seguenti attori principali:
\subsubsection{Client trasmittente}
Il dispositivo del client trasmette periodicamente via BLE un identificativo cifrato \(B_{\text{id}}\).
Questo valore non contiene informazioni personali o altri dati personali in chiaro.

\subsubsection{Client ricevente}
Il dispositivo del client ricevente ascolta i pacchetti BLE nell'area, li filtra per il nostro servizio, raccoglie i \(B_{\text{id}}\) e li invia al backend tramite
una richiesta autenticata. In cambio riceve i dettagli delle persone vicine.

\subsubsection{API e database}
Il backend riceve i \(B_{\text{id}}\) dai client riceventi, li risolve in utenti reali e restituisce i dettagli delle persone vicine.
\begin{itemize}
  \item conosce tutti gli utenti
  \item gestisce le chiavi crittografiche
  \item genera i pacchetti BLE \(B_{\text{id}}\) per i client trasmittenti
  \item risolve i \(B_{\text{id}}\) per i client riceventi
  \item comunica con il client ricevente tramite API le informazioni necessarie
\end{itemize}

\subsection{Principio di base}
Via radio BLE passa soltanto l'identificativo cifrato \(B_{\text{id}}\).
Tutto ciò che il client ricevente conosce sugli utenti vicini lo apprende soltanto tramite il backend via API, mai in chiaro su BLE.
In questo modo il canale radio resta molto povero di informazioni ed i dati sensibili rimangono confinati al backend.

\subsection{Ruolo del client}
L'app svolge tre funzioni principali:

\subsubsection{Advertising BLE}
Genera e trasmette periodicamente nei pacchetti BLE l'identificativo cifrato \(B_{\text{id}}\) valido per lo slot temporale corrente.
Ogni intervallo di tempo viene utilizzato un nuovo \(B_{\text{id}}\). Il client ha con sè in memoria i \(B_{\text{id}}\) generati per le successive
24 ore, che vengono mandati dal backend in un batch.

\subsubsection{Scanning BLE}
Ascolta i pacchetti BLE nell'area circostante, filtra quelli rilevanti per il nostro servizio e raccoglie i \(B_{\text{id}}\) visti.

\subsubsection{Risoluzione tramite backend}
Invia una richiesta al backend con i \(B_{\text{id}}\) raccolti, riceve in risposta i dettagli delle persone vicine e li mostra all'utente.\footnote{Sul dispositivo non esiste un mapping persistente \(B_{\text{id}} \rightarrow\) utente. La risoluzione di informazioni delle persone è sempre centralizzata nel backend.}

\subsection{Flusso logico}
Il flusso complessivo è il seguente:

\subsubsection{Generazione degli identificativi \(B_{\text{id}}\)}
Per ogni utente e per ogni intervallo di tempo il backend genera un valore cifrato \(B_{\text{id}}\) usando una cifratura AES-ENC ed una
chiave master AES-256. Il backend genera i successivi 144 \(B_{\text{id}}\) e li manda al client in un batch.

\subsubsection{Trasmissione via BLE}
Il client inserisce il \(B_{\text{id}}\) valido per lo slot temporale corrente nei pacchetti BLE e li trasmette periodicamente.

\subsubsection{Raccolta dei \(B_{\text{id}}\)}
Il client ricevente ascolta i pacchetti BLE nell'area, estrae i \(B_{\text{id}}\) e lo invia al backend per ottenere le informazioni dell'utente vicino.

\subsubsection{Risoluzione lato backend}
Il backend riceve i \(B_{\text{id}}\) dai client riceventi, e li decifra usando la chiave master AES-256 corretta. Dai \(B_{\text{id}}\) decifrati ottiene \(U_{\text{code}}\) ed il time slot.
Applica i controlli di validità temporale e restituisce i dettagli degli utenti vicini al client ricevente.

\subsubsection{Presentazione sull'APP}
Nell'app del client ricevente vengono mostrati i dettagli delle persone vicine ricevuti dal backend.

\section{Requisiti di progetto}

I requisiti principali del sistema sono:

\subsection{Privacy}
\begin{itemize}
  \item Nessun dato personale o sensibile deve essere trasmesso in chiaro via BLE;
  \item Un osservatore esterno che sniffa i pacchetti BLE non deve poter risalire in modo affidabile all'identità di un utente;
  \item Il mapping \(B_{\text{id}} \rightarrow\) utente deve essere noto solo al backend.
\end{itemize}


\subsection{Non tracciabilità}
\begin{itemize}
  \item Un osservatore esterno che sniffa i pacchetti BLE non deve poter seguire gli spostamenti di un utente nel tempo basandosi sull'identificativo \(B_{\text{id}}\);
  \item Anche disponendo di log storici dei \(B_{\text{id}}\) trasmessi non deve essere possibile risalire agli spostamenti di un utente.
\end{itemize}

\subsection{Efficienza BLE}
\begin{itemize}
  \item I pacchetti BLE devono essere il più piccoli possibile per ridurre il consumo energetico;
  \item Deve semplificare l'implementazione lato client tra gli stack BLE di iOS e Android.
\end{itemize}

\subsection{Scalabilità}
\begin{itemize}
  \item L'architettura deve poter supportare milioni di utenti attivi;
  \item La risoluzione di un singolo \(B_{\text{id}}\) deve essere \(\approx O(1)\), ovvero efficiente e non dipendente dal numero di utenti totali.
  \item Il backend deve poter gestire migliaia di richieste al secondo.
  \item Il client deve poter gestire centinaia di \(B_{\text{id}}\) raccolti.
\end{itemize}

\section{Modello dati dell'utente}

\textbf{Definizione (\(U_{\text{code}}\)):}
Sia \[U_{\text{code}} = \SI{64}{\bit} = \SI{8}{\byte}\] un identificativo univoco per ogni utente, generato casualmente al momento della registrazione.
Offre un namespace di \(2^{64} \approx \num{1.84e19}\) utenti, sufficiente per l'applicazione.
Non viene mai trasmesso in chiaro via BLE.

Ogni utente nel sistema è rappresentato dai seguenti dati:
\begin{itemize}
  \item identificativo univoco dell'utente \(U_{\text{code}}\);
  \item nome visualizzato;
\end{itemize}

Questi dati vengono usati solo nella comunicazione tra client ricevente e backend.

\section{Design crittografico}
\subsection{Obiettivo della cifratura}
L'obiettivo principale della cifratura è:
\begin{itemize}
  \item semplicità e efficienza di una soluzione AES con crittografia reversibile;
  \item autonomia del client nella rotazione degli identificativi \(B_{\text{id}}\) senza dover contattare il backend frequentemente;
\end{itemize}

\subsection{Finestra temporale}
Il tempo viene discretizzato in slot regolari di durata fissa.

\textbf{Definizione (\(\Delta t_{\text{slot}}\)):} Sia
\[
  \Delta t_{\text{slot}} = \SI{600}{\second} = \SI{10}{\minute}
\]
la durata di uno slot temporale.

\textbf{Definizione (\(t_{\text{unix}}\)):}
Sia
\[
  t_{\text{unix}} \in \mathbb{N}, \qquad [t_{\text{unix}}] = \si{\second},
\]
il tempo corrente in secondi trascorsi dal 1 gennaio 1970.

\textbf{Definizione (\(S_{\text{slot}}\)):}
Sia
\[
  S_{\text{slot}}
  = \left\lfloor \frac{t_{\text{unix}}}{\Delta t_{\text{slot}}} \right\rfloor
  \in \mathbb{N}
\]
l'indice dello slot temporale corrente.

\(S_{\text{slot}}\) è concettualmente pubblico e non segreto, in quanto può essere calcolato da
chiunque conoscendo l'ora corrente.

Quando un identificativo \(B_{\text{id}}\) viene decifrato nel backend, si ottiene in particolare un
indice di slot codificato
\[
  S_{\text{slot}}^{\text{dec}} \in \mathbb{N},
\]
che rappresenta lo slot temporale a cui il \(B_{\text{id}}\) è destinato.

\textbf{Definizione (\(\delta\)):}
Sia
\[
  \delta = \frac{\Delta t_{\text{slot}}}{5}
\]
la tolleranza temporale simmetrica per la risoluzione dei \(B_{\text{id}}\). Il suo valore è quindi
dinamicamente legato alla durata dello slot: se \(\Delta t_{\text{slot}}\) cambia, cambia
proporzionalmente anche \(\delta\).

Dallo slot decifrato si ricavano i tempi di inizio e fine dello slot codificato:
\[
  t_{\text{start}}^{\text{dec}}
  = S_{\text{slot}}^{\text{dec}} \cdot \Delta t_{\text{slot}}, \qquad
  t_{\text{end}}^{\text{dec}}
  = (S_{\text{slot}}^{\text{dec}} + 1) \cdot \Delta t_{\text{slot}}.
\]

Un identificativo \(B_{\text{id}}\) ricevuto al tempo \(t_{\text{unix}}\) è temporalmente valido se e
solo se
\[
  t_{\text{start}}^{\text{dec}} - \delta
  \;\le\;
  t_{\text{unix}}
  \;\le\;
  t_{\text{end}}^{\text{dec}} + \delta.
\]

In altre parole, il backend accetta un \(B_{\text{id}}\) se il tempo corrente cade all'interno dello
slot codificato, esteso di una tolleranza \(\delta\) prima dell'inizio e \(\delta\) dopo la fine.
Questa finestra allargata assorbe il drift di orologio tra client e backend e la latenza di rete,
senza allargare troppo l'esposizione temporale di un singolo \(B_{\text{id}}\).

\subsection{Chiavi AES e rotazione}
La generazione dei \(B_{\text{id}}\) si basa su chiavi master AES-256. Le chiavi non sono statiche e vengono ruotate periodicamente.


\textbf{Definizione (\(K_i\)):} Sia
\[
  K_i \in \{0,1\}^{256}, \quad |K_i| = 256 \,\text{bit} = 32 \,\text{byte}, \quad \forall i \in \mathbb{N}
\]

una chiave master AES-256 usata per generare i \(B_{\text{id}}\).
Ogni $K_i$ è generata tramite CSPRNG e memorizzata in modo sicuro nel backend. Le chiavi non derivano da una root comune, ma sono valori casuali forti e indipendenti.
A ciascun $K_i$ è associato un intervallo di tempo.

\textbf{Definizione (\(T_{K_i}, T_{K_i,\mathrm{grazia}}\)):} Sia
\[
  T_{K_i} = \SI{72}{\hour}
\]
la durata di validità della chiave master \(K_i\). Rispettivamente, sia
\[
  T_{K_i,\mathrm{grazia}} = \frac{T_{K_i}}{5}
\]
la durata del periodo di grazia successiva alla scadenza di \(K_i\), durante la quale \(K_i\) è ancora accettata per la risoluzione dei \(B_{\text{id}}\).


\textbf{Definizione (\((K_i)_{i\in\mathbb{N}}\)):} Sia
\[
  (K_i)_{i\in\mathbb{N}}
\]
la sequenza di chiavi master AES-256 usate per generare i \(B_{\text{id}}\). Ogni chiave è valida per un intervallo di tempo fisso.

\textbf{Definizione (\(t_{K_i}^{\text{start}}, t_{K_i}^{\text{end}}, t_{K_i}^{\text{grazia}}\)):} Sia
\[
  t_{K_i}^{\text{start}} =
  \begin{cases}
    t_{\text{now}},           & i = 1, \\
    t_{K_{i-1}}^{\text{end}}, & i > 1,
  \end{cases}
  \qquad i \in \mathbb{N},\, i \ge 1.
\]
il tempo di inizio validità di \(K_i\). Sia inoltre
\[
  t_{K_i}^{\text{end}} = t_{K_i}^{\text{start}} + T_{K_i}, \qquad \forall i \in \mathbb{N}.
\]
il tempo di fine validità di \(K_i\). Infine, sia
\[
  t_{K_i}^{\text{grazia}} = t_{K_i}^{\text{end}} + T_{K_i,\mathrm{grazia}}, \qquad \forall i \in \mathbb{N}.
\]
il tempo di grazia, dove \(K_i\) non è più usata per generare nuovi \(B_{\text{id}}\), ma è ancora accettata ed utilizzata per la risoluzione.

\subsubsection{Politica di rotazione}
La rotazione delle chiavi è periodica, e l'intervallo tra l'attivazione di chiavi consecutive è costante
\[
  T_{K_i} = t_{K_i}^{\text{start}} - t_{K_{i-1}}^{\text{start}}, \qquad \forall i > 1.
\]

$K_i$ è la chiave corrente nel tempo \(t_{\text{now}}\) se e solo se:
\[
  t_{K_i}^{\text{start}} \leq t_{\text{now}} < t_{K_i}^{\text{end}}
\]


Mentre $K_{i-1}$ è la chiave precedente ancora accettata nel periodo di grazia se e solo se:
\[
  t_{K_{i-1}}^{\text{end}} = t_{K_i}^{\text{start}} < t_{\text{now}} \leq t_{K_{i-1}}^{\text{grazia}}
\]

La politica di rotazione è progettata in modo che, per ogni $t_{now}$,
\[
  \forall t \in \mathbb{R}, \quad
  \left|
  \left\{
  i \in \mathbb{N}
  \;\middle|\;
  t_{K_i}^{\text{start}} \le t \le t_{K_i}^{\text{grazia}}
  \right\}
  \right| \le 2.
\]

\subsubsection{Effetti di sicurezza}
Questa architettura garantisce che:
\begin{itemize}
  \item l'impatto temporale di una eventuale compromissione di chiave è limitato a \(T_{K_i} + T_{K_i,\mathrm{grazia}}\);
  \item un attaccante che compromette \(K_i\) non può risalire ai \(B_{\text{id}}\) generati con chiavi precedenti \(K_{i-1}\), \(i-1 < i\);
  \item la rotazione delle chiavi \(K_i\), il periodo di grazia \(T_{K_i,\mathrm{grazia}}\) e la tolleranza temporale \(\delta\) sono gestiti in modo centralizzato e configurabile solo nel backend;
  \item tutto il materiale crittografico resta confinato nel backend.
\end{itemize}

\section{Flusso del client}
Descriviamo in modo più formale cosa fa il client, sia quando trasmette sia quando ascolta i pacchetti BLE \(B_{\text{id}}\)

\textbf{Definizione (\(T_{\text{batch}}\)):} Sia
\[
  T_{\text{batch}} = \SI{24}{\hour}
\]
la durata di validità di un batch di \(B_{\text{id}}\) inviato dal backend al client trasmittente.
Rappresenta l'intevallo di tempo in cui il client è in grado di generare i \(B_{\text{id}}\) localmente senza dover contattare il backend (offline).

\textbf{Definizione (\(N_{\text{batch}}\)):} Sia
\[
  N_{\text{batch}} = \frac{T_{\text{batch}}}{\Delta t_{\text{slot}}}
\]
il numero di \(B_{\text{id}}\) contenuti in un batch inviato dal backend al client trasmittente.

\subsection{Gestione advertising BLE}
\textbf{Definizione (\(\bigl(B_{\text{id},k}\bigr)_{k=1}^{N_{\text{batch}}}\)):} Sia
\[
  \mathbf{B}
  = \bigl(B_{\text{id},k}\bigr)_{k=1}^{N_{\text{batch}}},
  \qquad B_{\text{id},k} \in \{0,1\}^{128},\ \forall k \in \{1,\dots,N_{\text{batch}}\}.
\]
una sequenza ordinata di identificativi $B_{\text{id}}$

Il client trasmittente mantiene in memoria locale un batch di identificativi \(B_{\text{id}}\) fornito dal backend, valida per $N_{\text{batch}}$ di $S_{slot}$ consecutivi.

\textbf{Definizione (\(S_j\)):} Sia l'indice del primo slot del batch \(S_1 \in \mathbb{N}\) con
\[
  S_j = S_1 + (j-1), \qquad \forall j \in \{1,\dots,N_{\text{batch}}\},
\]
l'indice dello slot temporale coperto dal batch a cui è destinato l'identificativo \(B_{\text{id},j}\).

Nell'istante di advertising BLE, il client calcola lo slot corrente
\[
  S_{\text{slot}}^{\text{curr}}
  = \left\lfloor \frac{t_{\text{unix}}}{\Delta t_{\text{slot}}} \right\rfloor.
\]

\[
  j^\ast = S_{\text{slot}}^{\text{curr}} - S_1 + 1,
\]
ed usa l'identificativo \(B_{\text{id},j^\ast}\) da mandare a condizione che
\[
  1 \le j^\ast \le N_{\text{batch}}.
\]

Il ciclo operativo si può riassumere così:
\begin{enumerate}
  \item Il backend genera un batch di identificativi
        \[
          \mathbf{B} = \bigl(B_{\text{id},j}\bigr)_{j=1}^{N_{\text{batch}}}
        \]
        e l'indice del primo slot del batch \(S_1\), e li invia al client trasmittente.

  \item Il client memorizza localmente \(\mathbf{B}\) e \(S_1\), validi per \(N_{\text{batch}}\) slot temporali consecutivi.

  \item Ogni volta che deve trasmettere un pacchetto BLE, il client misura il tempo corrente \(t_{\text{unix}}\) e calcola
        \[
          S_{\text{slot}}^{\text{curr}}
          = \left\lfloor \frac{t_{\text{unix}}}{\Delta t_{\text{slot}}} \right\rfloor.
        \]

  \item Dal valore di \(S_{\text{slot}}^{\text{curr}}\) ricava l’indice nel batch
        \[
          j^\ast = S_{\text{slot}}^{\text{curr}} - S_1 + 1.
        \]

  \item Se
        \[
          1 \le j^\ast \le N_{\text{batch}},
        \]
        il client inserisce \(B_{\text{id},j^\ast}\) nel pacchetto BLE e lo trasmette.

  \item Se invece \(j^\ast < 1\) oppure \(j^\ast > N_{\text{batch}}\), il batch corrente non copre più lo slot attuale e il client deve ottenere un nuovo batch dal backend eliminando quello attuale.
\end{enumerate}

\subsection{Gestione scanning BLE}
Il client ricevente si occupa di ascoltare e registrare i $B_{\text{id}}$ degli altri dispositivi. Per ogni istante

La logica, per ogni pacchetto BLE ricevuto, è:
\begin{enumerate}
  \item filtrare solo i pacchetti BLE del nostro servizio ed ignorare gli altri;
  \item estrarre il $B_{\text{id}}$ dal pacchetto;
  \item inviare al backend la richiesta di risoluzione del pacchetto $B_{\text{id}}$
\end{enumerate}
Ogni $B_{\text{id}}$ viene quindi risolto al massimo una volta.

\subsection{Gestione batch e condizioni di errore}
Il client deve gestire alcuni casi particolari legati alla disponibilità di un batch valido e alla risposta del backend.

Un batch
\[
  \mathbf{B}
  = \bigl(B_{\text{id},j}\bigr)_{j=1}^{N_{\text{batch}}},
  \qquad B_{\text{id},j} \in \{0,1\}^{128},
\]
con primo slot \(S_1 \in \mathbb{N}\) è \emph{valido} all'istante \(t_{\text{unix}}\) se e solo se
\[
  S_{\text{slot}}^{\text{curr}}
  = \left\lfloor \frac{t_{\text{unix}}}{\Delta t_{\text{slot}}} \right\rfloor
\]
soddisfa
\[
  S_1 \le S_{\text{slot}}^{\text{curr}} \le S_1 + N_{\text{batch}} - 1.
\]
Equivalentemente, esiste un indice
\[
  j^\ast = S_{\text{slot}}^{\text{curr}} - S_1 + 1
\]
tale che
\[
  1 \le j^\ast \le N_{\text{batch}},
\]
e l'identificativo da usare è \(B_{\text{id},j^\ast}\).

\subsubsection{Primo avvio o nuovo login}
Al primo avvio il client non dispone di un batch valido, richiede immediatamente al backend un nuovo batch che copre le prossime $T_{\text{batch}}$

\subsubsection{Batch scaduto}
Se il batch corrente è scaduto, cioè
\[
  S_{\text{slot}}^{\text{curr}} > S_1 + N_{\text{batch}} - 1
\]

ed il client non riesce ad ottenere un nuovo batch dal backend (assenza di connettività o errore), allora:
\begin{itemize}
  \item il client sospende l'advertising BLE;
  \item riprenderà l'advertising solo dopo aver ricevuto un nuovo batch valido, per evitare di trasmettere $B_{\text{id}}$ non più risolvibili dal backend
\end{itemize}

In condizioni normali, finché esiste un batch valido e l'orologio del dispositivo resta sufficientemente sincronizzato, l'app può operare in modo autonomo per una durata pari a $T_{\text{batch}}$, riallineandosi al backend non appena necessario.

\section{Flusso di risoluzione lato server}
\subsection{Raccolta del \(B_{\text{id}}\) dal client}
Quando il client ha raccolto un $B_{\text{id}}$, invia una richiesta al backend.
Formalmente, una richiesta tipica contiene:
\begin{itemize}
  \item il pacchetto BLE \(B_{\text{id}}\) da risolvere
\end{itemize}

Una singola chiamata permette di risolvere un solo $B_{\text{id}}$

\subsection{Decifratura, selezione chiave e condizione di accettazione}
\textbf{Definizione (\(K_{\text{curr}}(t_{\text{unix}})\)):} Per ogni istante \(t_{\text{unix}}\) esiste al più un indice \(i \in \mathbb{N}\) tale che
\[
  t_{K_i}^{\text{start}} \le t_{\text{unix}} < t_{K_i}^{\text{end}}.
\]
Quando tale indice esiste, lo denotiamo con \(i^\ast\) e definiamo
\[
  K_{\text{curr}}(t_{\text{unix}}) := K_{i^\ast}.
\]

\textbf{Definizione (\(K_{\text{prev}}(t_{\text{unix}})\)):} Dato \(i^\ast\) come sopra, se \(i^\ast > 1\) e
\[
  t_{K_{i^\ast-1}}^{\text{end}} = t_{K_{i^\ast}}^{\text{start}}
  \quad\text{e}\quad
  t_{K_{i^\ast}}^{\text{start}} < t_{\text{unix}} \le t_{K_{i^\ast-1}}^{\text{grazia}},
\]
allora definiamo
\[
  K_{\text{prev}}(t_{\text{unix}}) := K_{i^\ast-1}.
\]

\textbf{Definizione (\(\mathcal{K}_{\mathrm{cand}}(t_{\text{unix}})\)):} Sia
\[
  \mathcal{K}_{\mathrm{cand}}(t_{\text{unix}})
  = \left\{
  K_i \;\middle|\;
  i \in \mathbb{N},\;
  t_{K_i}^{\text{start}} \le t_{\text{unix}} \le t_{K_i}^{\text{grazia}}
  \right\}.
\]
l'insieme delle chiavi candidate al tempo \(t_{\text{unix}}\)

Quindi le chiavi a disposizione per il backend sono:
\[
  \mathcal{K}_{\mathrm{cand}}(t_{\text{unix}}) =
  \begin{cases}
    \{K_{\text{curr}}(t_{\text{unix}})\},                                   &
    \text{se \(K_{\text{prev}}(t_{\text{unix}})\) non è definita},            \\[4pt]
    \{K_{\text{curr}}(t_{\text{unix}}), K_{\text{prev}}(t_{\text{unix}})\}, &
    \text{altrimenti}.
  \end{cases}
\]

Riassumendo il backend dispone di:
\begin{itemize}
  \item la chiave corrente sempre utilizzabile \(K_{\text{curr}}(t_{\text{unix}})\);
  \item la chiave precedente, utilizzabile se ancora nel suo periodo di grazia \(K_{\text{prev}}(t_{\text{unix}})\).
\end{itemize}

Ogni \(B_{\text{id}}\) è il risultato della cifratura AES-ENC di una coppia
\((U_{\text{code}}, S_{\text{slot}}^{(64)})\), dove \(S_{\text{slot}}^{(64)}\) è la
rappresentazione a 64 bit dell'indice di slot temporale.

Dal valore \(S_{\text{slot}}^{(64)}\) decifrato il backend ricava l'indice di slot
\[
  S_{\text{slot}}^{\text{dec}} \in \mathbb{N}.
\]

Al tempo di risoluzione \(t_{\text{unix}}\), il backend calcola lo slot corrente
\[
  S_{\text{slot}}^{\text{curr}}
  = \left\lfloor \frac{t_{\text{unix}}}{\Delta t_{\text{slot}}} \right\rfloor.
\]


\subsubsection{Algoritmo di risoluzione per un singolo $B_{\text{id}}$}
Dato un \(B_{\text{id}}\) ricevuto in un istante \(t_{\text{unix}}\), il backend esegue i passi seguenti.

\paragraph{Scelta delle chiavi da provare}
In base all'istante corrente \(t_{\text{unix}}\), il backend determina l'insieme delle chiavi candidate
\[
  \mathcal{K}_{\mathrm{cand}}(t_{\text{unix}})
  = \left\{
  K_i \;\middle|\;
  i \in \mathbb{N},\;
  t_{K_i}^{\text{start}} \le t_{\text{unix}} \le t_{K_i}^{\text{grazia}}
  \right\}.
\]

Per costruzione della politica di rotazione, vale
\[
  1 \le \bigl|\mathcal{K}_{\mathrm{cand}}(t_{\text{unix}})\bigr| \le 2,
\]
e le chiavi vengono considerate in ordine di priorità:
\[
  K_{\text{curr}}(t_{\text{unix}}) \quad\text{prima}, \qquad
  K_{\text{prev}}(t_{\text{unix}}) \quad\text{poi (se definita)}.
\]

\paragraph{Tentativi di decifratura}
Per ciascuna chiave \(K \in \mathcal{K}_{\mathrm{cand}}(t_{\text{unix}})\), il backend tenta la decifratura
\[
  \bigl(U_{\text{code}}^{(K)},\, S_{\text{slot}}^{(64),(K)}\bigr)
  = \mathrm{AES\text{-}DEC}_{K}\bigl(B_{\text{id}}\bigr),
\]
dove \(U_{\text{code}}^{(K)} \in \{0,1\}^{64}\) e \(S_{\text{slot}}^{(64),(K)} \in \{0,1\}^{64}\).

Se la decifratura fallisce (blocco non valido, padding errato o formato non coerente), la chiave \(K\) viene scartata per quel \(B_{\text{id}}\).
Tutte le coppie \(\bigl(U_{\text{code}}^{(K)}, S_{\text{slot}}^{(64),(K)}\bigr)\) ottenute da decifrature riuscite vengono trattate come
\emph{candidate} e passano ai passi di verifica successivi.

Se nessuna chiave in \(\mathcal{K}_{\mathrm{cand}}(t_{\text{unix}})\) produce una decifratura valida, il \(B_{\text{id}}\) viene scartato come invalido.

\paragraph{Verifica della finestra temporale}
Da ciascun valore \(S_{\text{slot}}^{(64),(K)}\) decifrato, il backend ricava l'indice di slot
\[
  S_{\text{slot}}^{\text{dec},(K)} \in \mathbb{N}.
\]

Dallo slot decifrato costruisce i tempi di inizio e fine dello slot codificato:
\[
  t_{\text{start}}^{\text{dec},(K)}
  = S_{\text{slot}}^{\text{dec},(K)} \cdot \Delta t_{\text{slot}}, \qquad
  t_{\text{end}}^{\text{dec},(K)}
  = \bigl(S_{\text{slot}}^{\text{dec},(K)} + 1\bigr) \cdot \Delta t_{\text{slot}}.
\]

All'istante di risoluzione \(t_{\text{unix}}\), la coppia decifrata con chiave \(K\) è
temporalmente valida se e solo se
\[
  t_{\text{start}}^{\text{dec},(K)} - \delta
  \;\le\;
  t_{\text{unix}}
  \;\le\;
  t_{\text{end}}^{\text{dec},(K)} + \delta.
\]

Solo le coppie \(\bigl(U_{\text{code}}^{(K)}, S_{\text{slot}}^{\text{dec},(K)}\bigr)\) che soddisfano
la condizione di finestra temporale vengono considerate per la fase di verifica utente. Se nessuna
delle decifrature valide soddisfa la finestra temporale, il backend rifiuta il
\(B_{\text{id}}\) come scaduto o troppo distante nel tempo.


\paragraph{Verifica utente}
Per ciascuna coppia temporalmente valida \(\bigl(U_{\text{code}}^{(K)}, S_{\text{slot}}^{\text{dec},(K)}\bigr)\), il backend esegue un lookup sul proprio database
per verificare l'esistenza di un utente associato a \(U_{\text{code}}^{(K)}\).

In forma concettuale, il backend cerca \(U_{\text{code}}^{(K)}\) tale che \(U_{\text{code}}^{(K)}\) corrisponda a un utente registrato;

\paragraph{Risultato}
Se esiste esattamente una decifratura che supera sia la verifica temporale sia la verifica utente, allora il \(B_{\text{id}}\) viene considerato
risolto e il backend restituisce al client ricevente i dettagli dell'utente associato.

Se nessuna decifratura supera entrambe le verifiche, il \(B_{\text{id}}\) viene rifiutato. Per costruzione del modello (spazio delle chiavi e
disegno del formato), la probabilità che più decifrature diverse producano coppie coerenti e riferite a utenti validi è:

\[
  \Pr[\text{collisione di decifratura valida}] \le \varepsilon,
  \qquad \varepsilon \approx 2^{-128}.
\]

\paragraph{Sintesi}
In forma compatta, un \(B_{\text{id}}\) ricevuto al tempo \(t_{\text{unix}}\) viene accettato se e solo se
esiste almeno una chiave \(K \in \mathcal{K}_{\mathrm{cand}}(t_{\text{unix}})\) tale che
\[
  \exists\, U_{\text{code}} \in \{0,1\}^{64},\;
  S_{\text{slot}}^{\text{dec}} \in \mathbb{N} \ \text{tali che}
\]
\[
  \mathrm{AES\text{-}DEC}_{K}\bigl(B_{\text{id}}\bigr)
  = \bigl(U_{\text{code}}, S_{\text{slot}}^{(64)}\bigr),
\]
da cui si ottiene
\[
  S_{\text{slot}}^{\text{dec}} \in \mathbb{N}, \qquad
  t_{\text{start}}^{\text{dec}} =
  S_{\text{slot}}^{\text{dec}} \cdot \Delta t_{\text{slot}}, \quad
  t_{\text{end}}^{\text{dec}} =
  (S_{\text{slot}}^{\text{dec}} + 1) \cdot \Delta t_{\text{slot}}.
\]
La condizione temporale richiede che
\[
  t_{\text{start}}^{\text{dec}} - \delta
  \;\le\;
  t_{\text{unix}}
  \;\le\;
  t_{\text{end}}^{\text{dec}} + \delta,
\]
\[
  \text{e che \(U_{\text{code}\!}\) identifichi un utente valido/attivo nel backend al tempo } t_{\text{unix}}.
\]

Se nessuna chiave candidata soddisfa tutte le condizioni sopra, il \(B_{\text{id}}\) viene rifiutato.

\subsubsection{Risposta al client}
Per ogni \(B_{\text{id}}\) decifrato, il backend restituisce un risultato sintetico, senza esporre dettagli interni sulla causa di eventuali rifiuti.
In generale:
\begin{itemize}
  \item se tutti i controlli sono superati, il backend restituisce:
        \subitem un nome visualizzato
  \item se uno o più controlli falliscono, il backend restituisce una risposta neutra o un errore generico, senza fornire indizi aggiuntivi.
\end{itemize}

L'app lato client usa solo queste risposte aggregate per comporre gli utenti vicino a lui, senza mai accedere direttamente ai dettagli del processo di risoluzione interno al backend.

\section{Threat model e mitigazioni}
\subsection{Sniffing BLE passivo}
\paragraph{Rischio}
Un attaccante ascolta passivamente il traffico BLE in un’area e prova a ricostruire quali utenti
sono presenti o a costruire uno storico dei loro passaggi nel tempo.

\paragraph{Mitigazioni}
\begin{itemize}
  \item i pacchetti BLE contengono solo \(B_{\text{id}} \in \{0,1\}^{128}\), dove
        \[
          B_{\text{id}} = \mathrm{AES\text{-}ENC}_{K_i}\bigl(U_{\text{code}} \,\|\, S_{\text{slot}}^{(64)}\bigr),
        \]
        senza nomi, \(U_{\text{code}}\) in chiaro o altri identificativi stabili;
  \item per uno stesso utente, il valore \(B_{\text{id}}\) cambia a ogni slot temporale
        \(\Delta t_{\text{slot}}\), quindi la sequenza
        \(\bigl(B_{\text{id}}(S_{\text{slot}})\bigr)_{S_{\text{slot}}}\) è una famiglia
        di valori pseudocasuali dipendenti da \(S_{\text{slot}}\);
  \item il mapping
        \[
          B_{\text{id}} \longrightarrow (U_{\text{code}}, S_{\text{slot}})
        \]
        è noto solo al backend e richiede una chiave \(K_i\) valida:
        per un osservatore che vede solo il traffico radio, invertire AES-256 senza chiave
        è computazionalmente infattibile.
\end{itemize}

\subsection{Tracciamento nel tempo}
\paragraph{Rischio}
Seguire il percorso di un utente nel tempo basandosi su un identificativo fisso o correlabile.

\paragraph{Mitigazioni}
\begin{itemize}
  \item per un utente con identificativo \(U_{\text{code}}\) fissato, i valori
        \[
          B_{\text{id}}(S_{\text{slot}}) =
          \mathrm{AES\text{-}ENC}_{K_i}\bigl(U_{\text{code}} \,\|\, S_{\text{slot}}^{(64)}\bigr)
        \]
        sono diversi per slot diversi \(S_{\text{slot}}\);
  \item senza conoscere \(K_i\) e \(U_{\text{code}}\), un attaccante che osserva
        \(B_{\text{id}}(S_{\text{slot}_1})\) e \(B_{\text{id}}(S_{\text{slot}_2})\) non ha
        un criterio praticabile per stabilire se provengono dallo stesso utente o da utenti diversi;
  \item la correlazione tra slot diversi è possibile solo nel backend, dove è noto sia
        \(U_{\text{code}}\) sia la chiave \(K_i\).
\end{itemize}

\subsection{Replay e spoofing di \(B_{\text{id}}\)}
\paragraph{Rischio}
Un attaccante ripete (replay) un \(B_{\text{id}}\) osservato in passato o in un luogo diverso per
simulare la presenza di un altro utente in un istante o in un’area diversa.

\paragraph{Mitigazioni}
\begin{itemize}
  \item ogni \(B_{\text{id}}\) incorpora uno slot temporale \(S_{\text{slot}}^{\text{dec}}\); al tempo
        \(t_{\text{unix}}\) il backend calcola lo slot corrente
        \[
          S_{\text{slot}}^{\text{curr}}
          = \left\lfloor \frac{t_{\text{unix}}}{\Delta t_{\text{slot}}} \right\rfloor;
        \]
  \item la condizione di finestra temporale impone che un \(B_{\text{id}}\) sia accettato solo se
        il tempo corrente \(t_{\text{unix}}\) cade nell'intervallo
        \[
          t_{\text{start}}^{\text{dec}} - \delta
          \;\le\;
          t_{\text{unix}}
          \;\le\;
          t_{\text{end}}^{\text{dec}} + \delta,
        \]
        dove
        \(t_{\text{start}}^{\text{dec}} =
        S_{\text{slot}}^{\text{dec}} \cdot \Delta t_{\text{slot}}\),
        \(t_{\text{end}}^{\text{dec}} =
        (S_{\text{slot}}^{\text{dec}} + 1) \cdot \Delta t_{\text{slot}}\)
        e \(\delta = \Delta t_{\text{slot}}/5\);
  \item un \(B_{\text{id}}\) osservato in passato smette di essere risolvibile non appena
        \(t_{\text{unix}} < t_{\text{start}}^{\text{dec}} - \delta\) oppure
        \(t_{\text{unix}} > t_{\text{end}}^{\text{dec}} + \delta\);

  \item il replay di \(B_{\text{id}}\) fuori dalla sua finestra di validità viene quindi rifiutato
        dal backend come temporalmente invalido.
\end{itemize}

\subsection{Compromissione di chiavi o sistemi backend}
\paragraph{Rischio}
Accesso non autorizzato alle chiavi AES \(K_i\) o ai sistemi che le gestiscono e le usano per
generare e risolvere i \(B_{\text{id}}\).

\paragraph{Mitigazioni}
\begin{itemize}
  \item le chiavi \(K_i \in \{0,1\}^{256}\) sono generate tramite CSPRNG, indipendenti tra loro e
        non derivate da una chiave root comune;
  \item ogni chiave \(K_i\) è valida per un intervallo temporale limitato
        \[
          T_{K_i},
          \qquad
          T_{K_i,\mathrm{grazia}} = \frac{T_{K_i}}{5},
        \]
        per un’esposizione complessiva massima \(T_{K_i} + T_{K_i,\mathrm{grazia}}\);
  \item la compromissione di una singola chiave \(K_i\) permette, nel caso peggiore, di risolvere
        solo i \(B_{\text{id}}\) generati con quella chiave e nel relativo intervallo temporale;
        non consente di recuperare chiavi \(K_j\) con \(j \neq i\);
  \item l’accesso alle chiavi è confinato al backend, e nessuna chiave viene mai distribuita ai client.
\end{itemize}

\subsection{Compromissione o furto del dispositivo utente}
\paragraph{Rischio}
Un attaccante ottiene accesso al dispositivo dell’utente e può leggere il batch locale di
identificativi \(B_{\text{id}}\) memorizzato per l’advertising.

\paragraph{Mitigazioni}
\begin{itemize}
  \item il batch locale
        \[
          \mathbf{B}
          = \bigl(B_{\text{id},j}\bigr)_{j=1}^{N_{\text{batch}}}
        \]
        copre al più un orizzonte temporale di durata
        \[
          T_{\text{batch}},
          \qquad
          N_{\text{batch}} = \frac{T_{\text{batch}}}{\Delta t_{\text{slot}}},
        \]
        quindi l’attaccante può riutilizzare (replay) questi \(B_{\text{id},j}\) solo entro
        quella finestra temporale;
  \item i \(B_{\text{id},j}\) sono comunque valori cifrati
        \(\in \{0,1\}^{128}\): senza le chiavi \(K_i\), l’attaccante non può risalire a
        \(U_{\text{code}}\) né agli slot \((S_{\text{slot}}^{\text{dec}})\) in chiaro;
  \item le chiavi AES \(K_i\) non risiedono mai sul dispositivo: la compromissione di un singolo
        telefono non compromette il resto del sistema, né permette di derivare chiavi o
        identificativi futuri;
  \item l’impatto della compromissione di un device è quindi limitato alla possibilità di replay
        dei \(B_{\text{id}}\) contenuti nel batch corrente, per una durata massima pari a
        \(T_{\text{batch}}\), ed è considerato accettabile rispetto al beneficio di autonomia del client.
\end{itemize}

\section{Alternative considerate e modelli mitigati}
\subsection{Nome in chiaro via BLE}

\textbf{Idea}
Trasmettere direttamente un campo testuale \texttt{name} nei pacchetti
BLE come identificativo dell'utente.

\textbf{Motivi di scarto}

\begin{itemize}[nosep]
  \item esporrebbe un dato personale direttamente in chiaro a chiunque nel raggio BLE, cioè
        \[
          \texttt{name} \longrightarrow \text{persona reale}
        \]
        tramite un mapping banale e stabile;
  \item renderebbe immediato il tracciamento nel tempo: lo stesso \texttt{name} verrebbe osservato
        in luoghi e istanti diversi, permettendo di costruire una traiettoria temporale per utente;
  \item soffrirebbe dei limiti di payload BLE per nomi lunghi o caratteri Unicode complessi;
  \item è in contrasto con il principio di minimizzazione dei dati: il canale radio veicolerebbe
        più informazione del necessario rispetto all’obiettivo (\(B_{\text{id}}\) opaco è sufficiente).
\end{itemize}

\subsection{Identificativo statico non rotante}

\textbf{Idea}

Assegnare a ogni utente un identificativo fisso
\[
  U_{\text{static}} \in \{0,1\}^{128}
\]
e trasmettere sempre e solo \(U_{\text{static}}\) nei pacchetti BLE.

\textbf{Motivi di scarto}

\begin{itemize}[nosep]
  \item permetterebbe tracciamento perfetto: per uno stesso utente la sequenza
        \(\bigl(U_{\text{static}}(t_k)\bigr)_k\) è costante nel tempo, quindi un attaccante può
        costruire una funzione
        \[
          t \longmapsto \text{posizione}(U_{\text{static}}, t)
        \]
        a partire dagli sniffing;
  \item anche se \(U_{\text{static}}\) è pseudonimo, è un identificatore permanente e linkabile:
        una volta correlato a un utente reale, tutte le osservazioni passate e future diventano
        immediatamente attribuibili;
  \item non rispetta il requisito di non tracciabilità tra slot temporali.
\end{itemize}

\subsection{Identificativi random rotanti con mapping esplicito}

\textbf{Idea}

Per ogni utente e per ogni slot \(S_{\text{slot}} \in \mathbb{N}\), il backend genera un token casuale
\[
  B_{\text{id}}^{\mathrm{rand}} \gets \text{Random}_{128\text{-bit}}() \in \{0,1\}^{128}
\]
e salva una riga in una tabella di mapping esplicita
\[
  B_{\text{id}}^{\mathrm{rand}} \longrightarrow
  \bigl(U_{\text{code}}, S_{\text{slot}}\bigr).
\]

Il client trasmette \(B_{\text{id}}^{\mathrm{rand}}\) via BLE; alla risoluzione il backend esegue una
lookup su questa tabella.

\textbf{Motivi di scarto}

\begin{itemize}[nosep]
  \item il mapping esplicito
        \[
          B_{\text{id}}^{\mathrm{rand}} \longrightarrow (U_{\text{code}}, S_{\text{slot}})
        \]
        esiste nel database: se la tabella viene compromessa, un attaccante che ha sniffato
        \(B_{\text{id}}^{\mathrm{rand}}\) nel tempo può ricostruire in chiaro coppie
        \((U_{\text{code}}, S_{\text{slot}})\);
  \item richiede di generare e \emph{persistere} un nuovo token per ogni utente e per ogni slot:
        lo spazio cresce come
        \[
          O\bigl(|\mathcal{U}| \cdot N_{\text{slot}}\bigr),
        \]
        dove \(|\mathcal{U}|\) è il numero di utenti e \(N_{\text{slot}}\) il numero di slot coperti
        dallo storico; ciò implica tabelle molto grandi e job di purge frequenti;
  \item a parità di comportamento radio (identificativi opachi e rotanti), introduce più stato e
        complessità rispetto alla soluzione AES corrente, in cui il mapping è \emph{implicito}
        nel blocco cifrato:
        \[
          B_{\text{id}} = \mathrm{AES\text{-}ENC}_{K_i}\bigl(U_{\text{code}} \,\|\, S_{\text{slot}}^{(64)}\bigr),
        \]
        e viene ricostruito al volo con una sola AES-DEC.
\end{itemize}

\subsection{Modello HMAC-SHA-256 con tabelle di lookup}

\textbf{Idea}

Per ogni utente si definisce una chiave HMAC
\[
  k_{\text{HMAC}} \in \{0,1\}^{\ell}
\]
e per ogni slot \(S_{\text{slot}}\) si calcola un token
\[
  T = \text{HMAC-SHA-256}\bigl(k_{\text{HMAC}},\, U_{\text{code}} \,\|\, S_{\text{slot}}^{(64)}\bigr),
\]
che viene poi troncato a 16 byte per poter essere inserito nei pacchetti BLE:
\[
  B_{\text{id}}^{\mathrm{HMAC}} = \text{Trunc}_{128}(T) \in \{0,1\}^{128}.
\]

Il backend pre-calcola, per ogni utente e slot, i token \(B_{\text{id}}^{\mathrm{HMAC}}\) e mantiene
una tabella di mapping
\[
  B_{\text{id}}^{\mathrm{HMAC}} \longrightarrow (U_{\text{code}}, S_{\text{slot}}).
\]

\textbf{Vantaggi}

\begin{itemize}[nosep]
  \item se \(k_{\text{HMAC}}\) viene fornita al client, i \(B_{\text{id}}^{\mathrm{HMAC}}\) possono
        essere generati direttamente sul dispositivo a partire da \(U_{\text{code}}\) e
        \(S_{\text{slot}}\);
  \item la risoluzione lato backend diventa una lookup \(O(1)\) su tabella chiave–valore;
  \item HMAC-SHA-256 è considerato crittograficamente robusto (PRF) con chiave segreta.
\end{itemize}

\textbf{Svantaggi}

\begin{itemize}[nosep]
  \item richiede tabelle di lookup per slot (o per finestre di slot) da tenere in memoria o su
        storage molto veloce:
        \[
          O\bigl(|\mathcal{U}| \cdot N_{\text{slot}}\bigr) \text{ entries};
        \]
  \item impone job periodici che generano token per tutti gli utenti a ogni slot, con costo
        computazionale e I/O significativo al crescere di \(|\mathcal{U}|\);
  \item se \(k_{\text{HMAC}}\) viene distribuita al device, la compromissione del dispositivo
        espone la chiave per quell’utente, permettendo di generare \(B_{\text{id}}^{\mathrm{HMAC}}\)
        arbitrari per slot futuri;
  \item un database che contenga il mapping
        \[
          B_{\text{id}}^{\mathrm{HMAC}} \longrightarrow (U_{\text{code}}, S_{\text{slot}})
        \]
        e venga compromesso insieme a \(k_{\text{HMAC}}\) o ad altri segreti, permette la
        ricostruzione completa delle identità associate agli ID sniffati.
\end{itemize}

\textbf{Come è stato mitigato}

Nella soluzione attuale:

\begin{itemize}[nosep]
  \item manteniamo il vantaggio principale del modello HMAC (autonomia nella rotazione lato
        client), ma usiamo \(B_{\text{id}}\) generati in batch dal backend con AES:
        \[
          B_{\text{id}} = \mathrm{AES\text{-}ENC}_{K_i}\bigl(U_{\text{code}} \,\|\, S_{\text{slot}}^{(64)}\bigr);
        \]
  \item non distribuiamo chiavi ai device: il client riceve direttamente i \(B_{\text{id}}\)
        necessari per un orizzonte \(T_{\text{batch}} = \SI{24}{\hour}\), senza conoscere \(K_i\);
  \item evitiamo tabelle di lookup per slot: la risoluzione richiede solo un numero costante di
        decifrature AES-256 (una per ciascuna chiave candidata in
        \(\mathcal{K}_{\mathrm{cand}}(t_{\text{unix}})\)), sempre in \(O(1)\) rispetto al numero di utenti.
\end{itemize}

\subsection{Modello AES server-centrico precedente}

\textbf{Idea originaria}

Utilizzare una singola chiave master statica \(K_{\text{enc}} \in \{0,1\}^{256}\) lato backend per
cifrare
\[
  \texttt{plaintext\_block}
  = U_{\text{code}} \,\|\, S_{\text{slot}}^{(64)}.
\]

Per ogni slot \(S_{\text{slot}}\), il client richiede al backend il corrispondente
\[
  B_{\text{id}}(S_{\text{slot}})
  = \mathrm{AES\text{-}ENC}_{K_{\text{enc}}}
  \bigl(U_{\text{code}} \,\|\, S_{\text{slot}}^{(64)}\bigr)
\]
(e, al più, pochi slot futuri) e lo usa subito per l'advertising BLE. Non sono previsti batch
ampi, né rotazioni frequenti di chiave.

\textbf{Vantaggi}

\begin{itemize}[nosep]
  \item logica estremamente semplice lato client: per ogni slot, 1 richiesta \(\to\) 1 \(B_{\text{id}}\);
  \item risoluzione sempre in \(O(1)\) con una sola AES-DEC usando \(K_{\text{enc}}\);
  \item nessuna tabella di mapping \(B_{\text{id}} \to\) utente: il mapping è implicito nella
        decifratura.
\end{itemize}

\textbf{Svantaggi}

\begin{itemize}[nosep]
  \item dipendenza costante dalla rete per aggiornare il \(B_{\text{id}}\): in assenza di
        connettività il client non è in grado di avanzare negli slot e non può trasmettere valori
        validi;
  \item presenza di una chiave a lunga vita \(K_{\text{enc}}\): se compromessa, un attaccante può
        decifrare tutti i \(B_{\text{id}}\) passati e futuri e recuperare
        \((U_{\text{code}}, S_{\text{slot}})\) per ogni pacchetto sniffato;
  \item è difficile limitare il raggio d'impatto temporale di una chiave compromessa: l’esposizione
        non è limitata a una finestra \([t_{K_i}^{\text{start}}, t_{K_i}^{\text{grazia}}]\), ma
        potenzialmente all’intera storia del sistema.
\end{itemize}

\textbf{Come è stato mitigato}

La soluzione attuale può essere vista come una versione ``rafforzata'' del modello AES originario:

\begin{itemize}[nosep]
  \item mantiene la struttura base AES-256:
        \[
          B_{\text{id}} = \mathrm{AES\text{-}ENC}_{K_i}
          \bigl(U_{\text{code}} \,\|\, S_{\text{slot}}^{(64)}\bigr),
        \]
        e la risoluzione in \(O(1)\) tramite un numero costante di AES-DEC;
  \item introduce il batch di durata \(T_{\text{batch}} = \SI{24}{\hour}\) inviato al client, che riceve
        una sequenza
        \(\mathbf{B} = (B_{\text{id},j})_{j=1}^{N_{\text{batch}}}\) e può operare offline finché il batch
        resta valido;
  \item sostituisce la singola chiave \(K_{\text{enc}}\) con una sequenza di chiavi master
        \((K_i)_{i \in \mathbb{N}}\), ciascuna valida per un intervallo temporale limitato
        \[
          T_{K_i} = \SI{72}{\hour},
          \qquad
          T_{K_i,\mathrm{grazia}} = \frac{T_{K_i}}{5},
        \]
        riducendo l’impatto di un’eventuale compromissione a una finestra temporale finita;
  \item usa chiavi \(K_i\) generate in modo indipendente tramite CSPRNG, memorizzate e versionate
        in modo sicuro nel backend, senza mai distribuirle ai client.
\end{itemize}

\section{Sintesi complessiva}

In questa sezione riassumiamo in modo compatto il funzionamento complessivo del sistema, le
proprietà desiderate e i principali trade-off progettuali.

\subsection{Architettura in sintesi}

\begin{itemize}
  \item Ogni utente è identificato in modo univoco da
        \[
          U_{\text{code}} \in \{0,1\}^{64},
        \]
        generato casualmente al momento della registrazione e usato solo nelle comunicazioni
        sicure tra client e backend.

  \item Il tempo è discretizzato in slot di durata fissata
        \[
          \Delta t_{\text{slot}} = \SI{600}{\second} = \SI{10}{\minute},
        \]
        con indice di slot
        \[
          S_{\text{slot}} = \left\lfloor \frac{t_{\text{unix}}}{\Delta t_{\text{slot}}} \right\rfloor
          \in \mathbb{N}.
        \]

  \item Il backend mantiene una sequenza di chiavi master AES-256
        \[
          (K_i)_{i\in\mathbb{N}}, \qquad K_i \in \{0,1\}^{256},
        \]
        ciascuna valida per una finestra temporale
        \[
          [t_{K_i}^{\text{start}},\, t_{K_i}^{\text{end}}],
          \qquad T_{K_i} = t_{K_i}^{\text{end}} - t_{K_i}^{\text{start}},
        \]
        seguita da un periodo di grazia
        \[
          [t_{K_i}^{\text{end}},\, t_{K_i}^{\text{grazia}}],
          \qquad T_{K_i,\mathrm{grazia}} = \frac{T_{K_i}}{5}.
        \]
        La politica di rotazione garantisce che, per ogni istante \(t\), siano attive al più due
        chiavi:
        \[
          \left|
          \left\{
          i \in \mathbb{N}
          \;\middle|\;
          t_{K_i}^{\text{start}} \le t \le t_{K_i}^{\text{grazia}}
          \right\}
          \right| \le 2.
        \]

  \item Ogni identificativo trasmesso via BLE è un blocco AES cifrato:
        \[
          B_{\text{id}} =
          \mathrm{AES\text{-}ENC}_{K_i}\bigl(U_{\text{code}} \,\|\, S_{\text{slot}}^{(64)}\bigr)
          \in \{0,1\}^{128},
        \]
        dove \(S_{\text{slot}}^{(64)}\) è la rappresentazione a 64 bit dell’indice di slot.

  \item Per ogni utente, il backend genera periodicamente un batch
        \[
          \mathbf{B} = \bigl(B_{\text{id},j}\bigr)_{j=1}^{N_{\text{batch}}},
          \qquad N_{\text{batch}} = \frac{T_{\text{batch}}}{\Delta t_{\text{slot}}},
        \]
        associato a un primo slot \(S_1\), in modo che ogni elemento sia destinato a uno slot:
        \[
          S_j = S_1 + (j-1), \qquad B_{\text{id},j} \longrightarrow S_j.
        \]

  \item Il client trasmittente conserva localmente \(\mathbf{B}\) e \(S_1\) e, all’istante \(t\),
        calcola lo slot corrente
        \[
          S_{\text{slot}}^{\text{curr}}
          = \left\lfloor \frac{t_{\text{unix}}}{\Delta t_{\text{slot}}} \right\rfloor
        \]
        per ricavare l’indice
        \[
          j^\ast = S_{\text{slot}}^{\text{curr}} - S_1 + 1,
        \]
        usando \(B_{\text{id},j^\ast}\) se e solo se \(1 \le j^\ast \le N_{\text{batch}}\).

  \item Il client ricevente effettua solo scanning passivo dei pacchetti BLE, estrae i
        \(B_{\text{id}}\) relativi al servizio e li invia al backend per la risoluzione, senza
        possedere né chiavi né mapping locali persistenti.

  \item Alla risoluzione, il backend:
        \begin{enumerate}[nosep]
          \item determina l’insieme di chiavi candidate
                \(\mathcal{K}_{\mathrm{cand}}(t_{\text{unix}})\) (chiave corrente e, se in grazia,
                chiave precedente);
          \item tenta la decifratura di \(B_{\text{id}}\) con ciascuna \(K \in \mathcal{K}_{\mathrm{cand}}\);
          \item ricava \((U_{\text{code}}, S_{\text{slot}}^{\text{dec}})\) per le decifrature riuscite e
                calcola
                \[
                  t_{\text{start}}^{\text{dec}}
                  = S_{\text{slot}}^{\text{dec}} \cdot \Delta t_{\text{slot}}, \qquad
                  t_{\text{end}}^{\text{dec}}
                  = (S_{\text{slot}}^{\text{dec}} + 1) \cdot \Delta t_{\text{slot}};
                \]
          \item verifica la finestra temporale imponendo che
                \[
                  t_{\text{start}}^{\text{dec}} - \delta
                  \;\le\;
                  t_{\text{unix}}
                  \;\le\;
                  t_{\text{end}}^{\text{dec}} + \delta;
                \]

          \item verifica che \(U_{\text{code}}\) corrisponda a un utente valido/attivo;
          \item se esiste un’unica soluzione coerente, restituisce i dettagli dell’utente al client
                ricevente.
        \end{enumerate}
\end{itemize}

\subsection{Proprietà di sicurezza e privacy}

Il disegno complessivo soddisfa i requisiti principali:

\begin{itemize}
  \item \textbf{Minimizzazione sul canale radio}: via BLE transita solo
        \(B_{\text{id}} \in \{0,1\}^{128}\), ovvero un blocco cifrato. Nessun dato personale
        (nome, email, ecc.) o identificativo persistente viene trasmesso in chiaro.
  \item \textbf{Non tracciabilità a livello radio}: per un utente fisso con \(U_{\text{code}}\)
        costante, la sequenza
        \[
          \bigl(B_{\text{id}}(S_{\text{slot}})\bigr)_{S_{\text{slot}}}
          = \bigl(
          \mathrm{AES\text{-}ENC}_{K_i}(U_{\text{code}} \,\|\, S_{\text{slot}}^{(64)})
          \bigr)_{S_{\text{slot}}}
        \]
        produce valori diversi per slot diversi; un osservatore esterno non può, in pratica,
        collegare due \(B_{\text{id}}\) a uno stesso utente senza conoscere chiavi e dati interni.
  \item \textbf{Limitazione del replay}: ogni \(B_{\text{id}}\) è accettato solo se, dato lo slot
        decifrato \(S_{\text{slot}}^{\text{dec}}\) e i corrispondenti
        \[
          t_{\text{start}}^{\text{dec}}
          = S_{\text{slot}}^{\text{dec}} \cdot \Delta t_{\text{slot}}, \qquad
          t_{\text{end}}^{\text{dec}}
          = (S_{\text{slot}}^{\text{dec}} + 1) \cdot \Delta t_{\text{slot}},
        \]
        il tempo corrente \(t_{\text{unix}}\) soddisfa
        \[
          t_{\text{start}}^{\text{dec}} - \delta
          \;\le\;
          t_{\text{unix}}
          \;\le\;
          t_{\text{end}}^{\text{dec}} + \delta.
        \]
        Un replay al di fuori di questa finestra viene rifiutato come temporalmente invalido.

  \item \textbf{Impatto limitato di compromissione di chiave}: ogni chiave \(K_i\) governa un
        intervallo limitato \([t_{K_i}^{\text{start}}, t_{K_i}^{\text{grazia}}]\); la compromissione
        di \(K_i\) non permette di derivare né chiavi precedenti né successive, limitando il danno
        temporale a \(T_{K_i} + T_{K_i,\mathrm{grazia}}\).
  \item \textbf{Assenza di chiavi lato device}: i client non ricevono mai \(K_i\); conservano solo
        batch di \(B_{\text{id}}\) già cifrati. La compromissione di un device esposto consente
        al più il replay dei \(B_{\text{id}}\) del batch corrente, ma non la decifratura né la
        generazione di nuovi valori al di fuori dell’orizzonte \(T_{\text{batch}}\).
  \item \textbf{Mapping centralizzato e implicito}: il mapping
        \[
          B_{\text{id}} \longrightarrow (U_{\text{code}}, S_{\text{slot}})
        \]
        esiste solo implicitamente tramite la decifratura AES nel backend. Non esistono tabelle
        persistenti \((B_{\text{id}} \to U_{\text{code}})\) che, se compromesse, rivelerebbero
        direttamente le identità.
\end{itemize}

\subsection{Scalabilità e complessità}

Dal punto di vista computazionale:

\begin{itemize}
  \item la risoluzione di un singolo \(B_{\text{id}}\) richiede al più un numero costante di
        operazioni AES-DEC, pari a \(|\mathcal{K}_{\mathrm{cand}}(t_{\text{unix}})| \le 2\); la
        complessità è quindi \(\Theta(1)\) rispetto al numero totale di utenti;
  \item non sono necessarie tabelle di lookup proporzionali a
        \( |\mathcal{U}| \cdot N_{\text{slot}} \), come negli approcci puramente random o HMAC
        con pre-calcolo per slot;
  \item il carico sul backend cresce linearmente con il numero di \(B_{\text{id}}\) da risolvere
        nel tempo (numero di richieste), ma non con il numero di utenti registrati.
\end{itemize}

\subsection{Trade-off e limiti noti}

\begin{itemize}
  \item \textbf{Dipendenza dal backend per la risoluzione}: il client ricevente non può risolvere
        localmente i \(B_{\text{id}}\); è richiesta connettività verso il backend per ottenere i
        dettagli delle persone vicine.
  \item \textbf{Autonomia limitata lato advertising}: il client trasmittente è autonomo finché
        dispone di un batch valido. Se il batch scade e non è possibile contattare il backend,
        l’advertising viene sospeso per non emettere \(B_{\text{id}}\) non risolvibili.
  \item \textbf{Sincronizzazione temporale}: la correttezza della finestra temporale richiede che
        l’orologio del dispositivo non derivi eccessivamente nel tempo rispetto al backend; la
        tolleranza \(\delta\) assorbe drift e latenza, ma non copre derive arbitrarie.
  \item \textbf{Osservatore fisico}: come in tutti i modelli basati su prossimità radio, un
        osservatore fisico presente nello stesso ambiente può comunque inferire co-presenze
        (chi è vicino a chi) a livello di contesto reale, anche se non può invertire
        crittograficamente i \(B_{\text{id}}\).
\end{itemize}

Nel complesso, la soluzione proposta combina:

\begin{itemize}
  \item identificativi BLE opachi e rotanti basati su AES-256;
  \item rotazione periodica delle chiavi master con periodo finito e grace period;
  \item batch lato client per autonomia offline limitata e controllata;
  \item risoluzione centralizzata in \(\Theta(1)\) con un numero costante di decifrature;
  \item assenza di mapping espliciti persistenti tra \(B_{\text{id}}\) e utenti.
\end{itemize}

Questi elementi forniscono un compromesso bilanciato tra privacy, non tracciabilità, efficienza
BLE, scalabilità backend e semplicità di implementazione lato client.


\end{document}